% Corrected French for printed p.179; departures are append-only below.
\SourcePage{179}

\LayerTarget{tohoku:eq:special-cover-H1-is-Z:p2}
\noindent la limite $\check H^1(X,H^1(F))=\mathbf Z$, ce qui achève la
démonstration.

\LayerTarget{tohoku:section:3.9}
\textbf{3.9. Critères d'acyclicité par la méthode des
recouvrements.}\textsuperscript{8)}
\LayerTarget{tohoku:setup:acyclicity-cover-families}
Soient $X$ un espace topologique, $\mathfrak S$ un ensemble non vide de parties
de $X$. Pour tout $A\in\mathfrak S$, on suppose donné un ensemble
$\mathfrak R(A)$ non vide de recouvrements $\mathbf R$ de $A$ par des
ensembles qui sont dans $\mathfrak S$ ainsi que leurs intersections finies. On
suppose que si $\mathbf R,\mathbf R'\in\mathfrak R(A)$ et
$B\in\mathbf R$, alors la trace $\mathbf R'_B$ de $\mathbf R'$ sur $B$
appartient à $\mathfrak R(B)$.
\LayerTarget{tohoku:conditions:acyclicity-cover-families}
On suppose de plus que l'on est dans l'une des trois conditions suivantes
(qui permettent d'écrire la suite spectrale de Leray-Cartan pour chacun des
recouvrements $\mathbf R\in\mathfrak R(A)$, avec $A\in\mathfrak S$): (i) les
$A\in\mathfrak S$ sont ouverts (ii) Les $A\in\mathfrak S$ sont fermés et les
recouvrements $\mathbf R\in\mathfrak R(A)$ sont finis (iii). Les
$A\in\mathfrak S$ sont fermés, $X$ est paracompact, les
$\mathbf R\in\mathfrak R(A)$ sont localement finis.

\LayerTarget{tohoku:theorem:3.9.1}
\textsc{Théorème 3.9.1.} \textit{Sous les conditions précédentes,
supposons donnés un faisceau abélien $F$ sur $X$, et un entier naturel
$n\geqq 0$. Supposons les conditions suivantes vérifiées :}

\LayerTarget{tohoku:condition:3.9.1:A-n}
\begingroup\itshape
$A(n)$ : $H^i(\mathbf R,F)=0$ pour $1\leqq i\leqq n$ et tout
$\mathbf R\in\mathfrak R(A)$, $A\in\mathfrak S$.\par
\endgroup

\LayerTarget{tohoku:condition:3.9.1:B-n-minus-1}
\begingroup\itshape\hangindent=2.5em\hangafter=1
$B(n-1)$ : Pour tout $A\in\mathfrak S$, tout $c^i\in H^i(A,F)$ \textup{(}avec
$1\leqq i\leqq n-1$\textup{)} il existe une partie finie $L$ de $\mathfrak R(A)$ telle
que, si $\mathbf R^L$ désigne le recouvrement ``intersection'' des
recouvrements $\mathbf R\in L$, la restriction de $c^i$ à tout ensemble
$B\in\mathbf R^L$ est nulle.\par
\endgroup

\LayerTarget{tohoku:conclusion:3.9.1}
\noindent\textit{Sous ces conditions, pour tout $A\in\mathfrak S$ on a
$H^i(A,F)=\{0\}$ pour $1\leqq i\leqq n-1$, et si $c^n\in H^n(A,F)$, alors
$c^n$ est nul si et seulement si on peut trouver une partie finie $L$ de
$\mathfrak R(A)$ telle que la restriction de $c^n$ à tout
$B\in\mathbf R^L$ soit nulle.}

Enonçons tout de suite les corollaires les plus intéressants :

\LayerTarget{tohoku:corollary:3.9.1:1}
\textsc{Corollaire 1.} \textit{Avec les notations du théorème \textup{3.9.1},
pour qu'on ait $H^i(A,F)=\{0\}$ pour tout $A\in\mathfrak S$ et
$1\leqq i\leqq n$, il faut et il suffit que les conditions $A(n)$ et $B(n)$
soient satisfaites.}

\LayerTarget{tohoku:proof:corollary-3.9.1-1:sufficiency}
En effet, la suffisance résulte aussitôt du théorème.
\LayerTarget{tohoku:proof:corollary-3.9.1-1:necessity}
Réciproquement, supposons $H^i(A,F)=\{0\}$ pour $A\in\mathfrak S$ et
$1\leqq i\leqq n$, alors $B(n)$ est trivialement vérifiée, vérifions $A(n)$ :
la suite spectrale de Leray pour le recouvrement $\mathbf R$ de $A$,
(théorème 3.8.1 et remarque de 3.8) aboutit à $H^*(A,F)$ et a pour terme
initial
$E_2^{p,q}=H^p(\mathbf R,H^q(F))$,
qui est nul si $1\leqq q\leqq n$ car $C(\mathbf R,H^q(F))=0$ pour ces valeurs
de $q$ (les intersections finies d'ensembles $\in\mathbf R$ appartenant à
$\mathfrak S$). On en conclut classiquement que
$H^i(A,F)=H^i(\mathbf R,F)$ pour $0\leqq i\leqq n$, d'où
$H^i(\mathbf R,F)=0$ pour ces $i>0$.

\LayerTarget{tohoku:corollary:3.9.1:2}
\textsc{Corollaire 2.} \textit{Supposons \textup{(}avec les notations du
théorème \textup{3.9.1)} que pour tout $A\in\mathfrak S$ et tout recouvrement ouvert de
$A$, on puisse trouver un recouvrement plus fin de la forme $\mathbf R^L$, où
$L$ est une partie finie de $\mathfrak R(A)$. Alors pour que
$H^i(A,F)=\{0\}$ pour tout $A\in\mathfrak S$ et $1\leqq i\leqq n$, il faut et
il suffit que $H^i(\mathbf R,F)=\{0\}$ pour tout
$\mathbf R\in\mathfrak R(A)$ \textup{(}$A\in\mathfrak S$\textup{)} et $1\leqq i\leqq n$.}

\LayerTarget{tohoku:proof:corollary-3.9.1-2}
En effet, la condition $B(n)$ est vérifiée (quel que soit $n>0$) en vertu du
lemme 3.8.2 ; il suffit donc d'appliquer le corollaire 1.

\TohokuContinuation{180}
